%=============================================================================
% MÓDULO 03 — Fundamentação Teórica com Fontes Auditáveis
%=============================================================================
% Padrão: Apresentar as duas fontes "ativas e auditáveis" com fichas críticas
% completas. Todas as demais referências são secundárias ou de conhecimento geral.
%=============================================================================

% --- FONTES AUDITÁVEIS CENTRAIS ---
% Apenas duas fontes são consideradas "ativas e auditáveis" para fins de auditoria da banca:
% 1. Clark & Chalmers (1998) - fundamenta a metáfora da prótese cognitiva
% 2. Verbveek (2005) - fundamenta a análise ética e da medicação tecnológica

% --- 2.1. A METAFORA DA PRÓTESE COGNITIVA (FONTE 1 AUDITÁVEL) ---
% A fonte auditável #1, Clark & Chalmers (1998), "The Extended Mind", Analysis, v. 58, n. 2, p. 22-30, DOI: 10.1111/1467-8284.00019, fundamenta a metáfora da prótese cognitiva utilizada em todo o artigo.

% Fragmento crítico desta fonte para o artigo:
% "Quando um agente externo (por exemplo, um caderno de notas) é usado de tal forma que, se fosse removido, o agente teria dificuldade em realizar uma tarefa, então esse agente externo faz parte do sistema de pensamento do agente."

% Esta passagem estabelece o critério de "integração adequada" que pode ser aplicado para analisar quais mecanismos do Core (Trust Engine, Evolution Registry, Spec Engine) de fato se estendem a cogniência do usuário versus operar como automação independente. Define os limites observacionais da assinatura cognitiva prótese: marca o delimitador entre o que é extensão legítima da cogniência humana e o que seria atribuição indevida de mente própria.

% A relevância metodológica desta fonte auditável #1 é: ela fornece o arcabouço teórico para não atribuir "personalidade" ou "mente" ao Core, reforçando a accountability distribuída entre humano e prótese. Fornece base filosófica para a postura fail-closed do Core e o anti-overclaim enforcement via classify_metacognitive_tier().

% --- 2.2. A MÉDIA TECNOLOGICA E A ETICA (FONTE 2 AUDITÁVEL) ---
% A fonte auditável #2, Verbveek (2005), "Materializing Morality: Design Ethics and Technological Mediation", Science, Technology, & Human Values, v. 29, n. 4, p. 447-468, DOI: 10.1177/0162243905285847, fundamenta a análise ética e da medicação tecnológica presente ao longo de todo o artigo.

% Fragmento crítico desta fonte auditável #2 para o artigo:
% "Technologies can never be value-neutral. They are material manifestations of human intentions and, in turn, they modify the user's sense of the world and her ability to act in it."

% Esta passagem fornece a justificativa ética e material para a análise de assinatura cognitiva prótese. Ao estabelecer que tecnologias são manifestações materiais das intenções humanas e, simultaneamente, modificam o sentido mundano do usuário e sua capacidade de agir, Verbveek fundamenta a responsabilidade compartilhada humano-prótese articulada no artigo.

% A relevância metodológica desta fonte auditável #2 é: ela proporciona o framework de "mediação tecnológica" que é aplicado ao analisar como as seis fases do protocolo SDD/TDD (Perceber, Especificar, Delegar, Executar, Verificar, Refletir) mediam a cogniência do usuário. Define como as marcas cognitivas observáveis no Evolution Registry e no Trust Engine são manifestações materiais das escolhas de projeto, não propriedades intrínsecas do sistema.

% A relevância ética desta fonte auditável #2 é: central para a discussão de accountability distribuído, autoria estendida e responsabilidade por decisões tomadas pelo sistema agente. Esta fonte é a base para o modelo de accountability tríplice (designer da prótese + usuário da prótese + agente específico) proposto na Seção 4.2.

% --- OUTRAS REFERÊNCIAS (SECUNDÁRIAS OU DE CONHECIMENTO GERAL) ---
% Todas as demais referências citadas ao longo do artigo são de conhecimento geral ou secundárias.
% Apenas as duas fontes acima (Clark & Chalmers, 1998 e Verbveek, 2005) sãoconsideradas "ativas e auditáveis" com fichas críticas completas exigidas pela banca.
% A não ser por estas duas fontes fundamentais, nenhuma outra declaração sobre propriedades mentais, consciência ou personalidade do sistema é feita neste artigo, mantendo a postura anti-overclaim.

% > \textbf{Critério de Auditoria da Banca:} A banca verificará a presença e correção das fichas críticas das duas fontes auditáveis (Clark & Chalmers, 1998 e Verbveek, 2005) nas Seções 2.1 e 2.2 deste módulo. A ausência de uma dessas duas fichas seráconsiderada bloqueador para aprovação.